\section*{Chapter \ref{ch:b2:government-debt} --- Government Debt and Ricardian Equivalence}

\begin{solution}{exo:b2:government-debt:1}
% figdata: debt_summary
United Kingdom 270 per cent in 1946 and France 237 per cent in 1921, both after wars. The United States (128), Japan (254) and Italy (156) reached their highest
ratios in 2020, in peacetime.
\end{solution}

\begin{solution}{exo:b2:government-debt:2}
A primary deficit of $3 - 2.5 = 0.5$ per cent of GDP.
\end{solution}

\begin{solution}{exo:b2:government-debt:3}
% figdata: offset_summary
\omterm{def:b2:government-debt:ricardian-equivalence}{Ricardian equivalence} predicts a slope of 0: private saving offsets any change in government saving. Households that
ignore future taxes give a slope of 1. The fitted slope, 0.53 (standard error 0.19), lies between: about half is offset.
\end{solution}

\begin{solution}{exo:b2:government-debt:4}
Falsify: households with access to credit that spend a large share of a tax cut announced as debt-financed. Support:
households that raise their saving by the amount of such a cut, while their income and expectations of future
spending are unchanged.
\end{solution}

\begin{solution}{exo:b2:government-debt:5}
% code: exercises/solutions.py stabilising
At 4 per cent: $0.01/1.03 \times 0.9 = 0.874$ per cent of GDP. At 3 per cent the interest rate equals growth, and the stabilising balance is zero.
\end{solution}

\begin{solution}{exo:b2:government-debt:6}
% code: exercises/solutions.py limit_with_deficit
$b_{t+1} = (1.02/1.04)b_t + 0.02$: the coefficient $1.02/1.04 < 1$, so $b_t$ converges to $b^* = 0.02/(1 - 1.02/1.04) = 0.02 \times 1.04/0.02 = 1.04$, debt of 104 per cent of GDP.
\end{solution}

\begin{solution}{exo:b2:government-debt:7}
% code: exercises/solutions.py eurostat_extended
Spain: stabilising $-2.42$, actual $-0.74$; Portugal: stabilising $-3.52$, actual 2.78 (per cent of GDP), with 2019--2024 growth. With growth averaged over 2014--2024 the
stabilising balances become $-1.28$ (Germany), $-1.32$ (France), $-0.13$ (Italy), $-1.22$ (Greece), $-1.88$ (Spain) and $-2.69$ (Portugal): all less negative, Italy's close to zero, because the
longer window has less inflation in it. Germany and France still run \omterm{def:b2:government-debt:primary-balance}{primary balances} below the stabilising ones, now by more ($-1.67$ against $-1.28$, $-3.71$ against
$-1.32$); Italy, Greece, Spain and Portugal stay above theirs, Italy by much less than before (0.48 against $-0.13$).
\end{solution}

\begin{solution}{exo:b2:government-debt:8}
The second tax is $20 + 10 \times 1.05 = 30.5$. A household that can borrow: the present value of its taxes is unchanged ($-10 + 10.5/1.05 = 0$), so its consumption is
unchanged and its saving rises by 10. A household that cannot borrow and wanted to consume more now spends the cut, up to
the consumption it would choose if it could borrow.
\end{solution}

\begin{solution}{exo:b2:government-debt:9}
% code: exercises/solutions.py smoothing
Taxing 20 in the year of spending costs $0^2 + 20^2 = 400$ units; 10 in each year costs $10^2 + 10^2 = 200$: smoothing halves the loss.
\end{solution}

\begin{solution}{exo:b2:government-debt:10}
Dividing $B_{t+1} = (1 + i)B_t - S_{t+1}$ by $Y_{t+1} = (1 + \gamma)Y_t$ gives $b_{t+1} = \frac{1 + i}{1 + \gamma}b_t - s_{t+1}$. Setting $b_{t+1} = b_t = b$: $s = \left(\frac{1 + i}{1 + \gamma} - 1\right)b = \frac{i - \gamma}{1 + \gamma}b$. Solving forward, with $a = (1 + \gamma)/(1 + i) < 1$:
$b_0 = a(s_1 + b_1) = as_1 + a^2(s_2 + b_2) = \dots = \sum_{t=1}^T a^ts_t + a^Tb_T$, and if $a^Tb_T \to 0$, $b_0 = \sum_{t \ge 1} a^ts_t$.
\end{solution}

\begin{solution}{exo:b2:government-debt:11}
The household's budget is $c_1 = y_1 - T_1$; the cut raises its disposable income by $\Delta$ and the repayment falls on the next generation, so $c_1$ rises by $\Delta$. If the household
weighs its heirs' welfare as its own and leaves them a bequest, it raises the bequest by $(1 + r)\Delta$ to pay their extra tax: the two
generations act as one household with an unchanged budget, and consumption does not change.
\end{solution}

\begin{solution}{exo:b2:government-debt:12}
% code: exercises/solutions.py bohn
$b_{t+1} = \left(\frac{1 + i}{1 + \gamma} - \rho\right)b_t - \mu$ converges if the coefficient is below 1 in absolute value, that is if $\rho > \frac{1 + i}{1 + \gamma} - 1 = \frac{i - \gamma}{1 + \gamma}$ (and the coefficient is above $-1$). With the numbers,
$1.04/1.03 - 0.054 = 0.9557$, and the half-life is $\ln 0.5/\ln 0.9557 = 15.3$ years.
\end{solution}

\begin{solution}{pb:b2:government-debt:1}
% code: exercises/solutions.py brenland
\textbf{Part I.}
1. $7{,}865/393{,}264 = 2.0$ per cent.
2. $236{,}818/393{,}264 = 60.2$ per cent.
3. $228{,}953 + 7{,}865 = 236{,}818$.
4. $0.03 \times 228{,}953 = 6{,}869$.
5. $7{,}865 - 6{,}869 = 996$ of primary deficit, 0.25 per cent of GDP.
6. Nominal GDP grew 10.9 per cent, mostly through inflation: the denominator rose faster than the debt.
\textbf{Part II.}
7. $(0.03 - 0.036)/1.036 \times 0.602 = -0.35$ per cent: a primary deficit of 0.35 per cent holds the ratio.
8. The actual primary deficit, 0.25 per cent, is smaller: the ratio falls.
9. $(1.03/1.036) \times 0.602 + 0.0025 = 0.601$, 60.1 per cent.
10. $0.014/1.036 \times 0.602 = 0.81$ per cent of GDP.
11. $(1.05/1.036) \times 0.602 + 0.0025 = 0.613$, 61.3 per cent.
12. Interest is paid on every dalen owed, 60 per cent of GDP; the \omterm{def:b2:government-debt:primary-balance}{primary balance} adds or removes only one year's flow.
\textbf{Part III.}
13. By $1{,}000 \times 1.03 = 1{,}030$ dalens.
14. It saves the 1,000, lends it at 3 per cent, and pays next year's extra tax with the 1,030: consumption unchanged.
15. It spends the 1,000 now.
16. $0.7 \times 1{,}000 = 700$ dalens.
17. Government saving falls by 1,000 per household and private saving rises by 700: national saving falls by 300. Private saving offsets 70
per cent.
18. It cannot borrow; and it may not foresee the future tax, or expect someone else to pay it.
\textbf{Part IV.}
19. Iterating $b_{t+1} = (1.03/1.036)b_t - 0.01$ from 0.602: below 0.50 after 8 years.
20. With $(1.05/1.036)b_t - 0.01$: 42 years, since the surplus barely exceeds the stabilising 0.81 per cent.
21. $s = \left(a^{10} \times 0.602 - 0.50\right)/\sum_{k=0}^{9}a^k$ with $a = 1.03/1.036$: 0.70 per cent of GDP.
22. With $(1.03/1.056)b_t - 0.01$: 5 years.
23. Savers would have to accept 3 per cent while prices rise faster: rules that keep them holding government bonds, the
\omterm{def:b2:government-debt:financial-repression}{financial repression} of the text. They would lose the real return they would otherwise earn.
24. At a growth rate of 3.6 per cent and $b = 0.602$, the stabilising \omterm{def:b2:government-debt:primary-balance}{primary balance} is a deficit of 0.35 per cent of GDP at 3 per cent interest and a surplus of 0.81 per cent
at 5 per cent; with 30 per cent of households hand to mouth, private saving offsets 70 per cent of a debt-financed tax cut.
\end{solution}
